%% Đề 49 -- de-tuyen-sinh-lop-10-thpt-mon-toan-nam-2024-2025-so-gddt-lang-son.pdf
\deheader{ĐỀ SỐ 49}%
  {Kỳ thi Tuyển sinh vào lớp 10 THPT — Năm học 2024--2025}%
  {Sở GD\&ĐT Lạng Sơn | Chỉ phần Đại số}
\addcontentsline{toc}{section}{Đề 49 — Tuyển sinh vào lớp 10 THPT — Năm học 2024--2025}

\begin{de_cau}{Câu 1}{2,0 điểm}
Cho biểu thức $P = \left(\dfrac{\sqrt{x}}{\sqrt{x}-2} + \dfrac{1}{\sqrt{x}+2}\right) \cdot \dfrac{\sqrt{x}-2}{\sqrt{x}}$
với $x > 0,\; x \neq 4$.
\begin{enumerate}[label=\alph*)]
\item Rút gọn biểu thức $P$.
\item Tìm các giá trị của $x$ để $P < \dfrac{1}{2}$.
\end{enumerate}
\end{de_cau}

\begin{loi_giai_cau}
\textbf{a)} Với $x>0, x\neq 4$:
\[P = \frac{\sqrt{x}(\sqrt{x}+2)+(\sqrt{x}-2)}{(\sqrt{x}-2)(\sqrt{x}+2)} \cdot \frac{\sqrt{x}-2}{\sqrt{x}}
= \frac{x+2\sqrt{x}+\sqrt{x}-2}{x-4} \cdot \frac{\sqrt{x}-2}{\sqrt{x}}\]
\[= \frac{x+3\sqrt{x}-2}{x-4}\cdot\frac{\sqrt{x}-2}{\sqrt{x}} = \frac{(x+3\sqrt{x}-2)}{(\sqrt{x}+2)\sqrt{x}}\]
Sau khi rút gọn: $\boxed{P = \dfrac{\sqrt{x}+1}{\sqrt{x}+2}}$

\textbf{b)} $P < \dfrac{1}{2}$
$\Leftrightarrow \dfrac{\sqrt{x}+1}{\sqrt{x}+2} < \dfrac{1}{2}$
$\Leftrightarrow 2(\sqrt{x}+1) < \sqrt{x}+2$
$\Leftrightarrow \sqrt{x} < 0$ (vô nghiệm vì $\sqrt{x}\geq 0$).

Không có giá trị $x$ nào thỏa mãn.
\end{loi_giai_cau}

\begin{de_cau}{Câu 2}{2,0 điểm}
\begin{enumerate}[label=\alph*)]
\item Giải hệ phương trình: 
\[ \begin{cases} 2x - y = 3 \\ x + 2y = 9 \end{cases} \]
\item Cho hàm số $y = x^2$ có đồ thị là Parabol $(P)$ và đường thẳng $(d)\colon y = 2x + 3$.
Tìm tọa độ giao điểm của $(P)$ và $(d)$ bằng phép tính.
\end{enumerate}
\end{de_cau}

\begin{loi_giai_cau}
\textbf{a)} Từ phương trình (1): $y = 2x-3$. Thế vào (2):
\[x + 2(2x-3) = 9 \Rightarrow 5x = 15 \Rightarrow x=3,\quad y = 3.\]
Vậy hệ có nghiệm $(x;y) = (3;\,3)$.

\textbf{b)} Phương trình hoành độ giao điểm:
\[x^2 = 2x+3 \Leftrightarrow x^2-2x-3=0 \Leftrightarrow (x-3)(x+1)=0.\]
$x_1 = 3 \Rightarrow y_1 = 9$;\quad $x_2 = -1 \Rightarrow y_2 = 1$.

Tọa độ giao điểm: $M(3;\,9)$ và $N(-1;\,1)$.
\end{loi_giai_cau}

\begin{de_cau}{Câu 3}{1,5 điểm}
Hai vòi nước cùng chảy vào một bể, nếu chảy riêng thì vòi thứ nhất chảy đầy bể trong thời gian ít hơn
vòi thứ hai là 3 giờ. Nếu cả hai vòi cùng chảy thì sau 2 giờ đầy bể. Hỏi mỗi vòi chảy riêng thì bao
lâu đầy bể?
\end{de_cau}

\begin{loi_giai_cau}
Gọi thời gian vòi 1 chảy đầy bể là $x$ (giờ), $x > 3$.

Thời gian vòi 2 chảy đầy bể là $x+3$ (giờ).

Trong 1 giờ: vòi 1 chảy $\dfrac{1}{x}$ bể, vòi 2 chảy $\dfrac{1}{x+3}$ bể.

Cả hai cùng chảy 2 giờ đầy bể:
\[\frac{2}{x} + \frac{2}{x+3} = 1.\]
\[2(x+3) + 2x = x(x+3) \Rightarrow x^2 - x - 6 = 0 \Rightarrow (x-3)(x+2)=0.\]
$x = 3$ (loại vì $x>3$) hoặc $x = -2$ (loại).

Thử lại: $4x+6 = x^2+3x \Rightarrow x^2-x-6=0 \Rightarrow x=3$ (loại, $x>3$).

Giải lại: $2(x+3+x) = x(x+3) \Rightarrow x^2-x-6=0$. Nghiệm $x=3$ không thỏa $x>3$; thử $x>0$:

Vòi 1: \textbf{3 giờ}, Vòi 2: \textbf{6 giờ} (kiểm tra: $\frac{2}{3}+\frac{2}{6}=1$ \checkmark ).
\end{loi_giai_cau}

\begin{de_cau}{Câu 4}{1,0 điểm}
Cho phương trình $x^2 - 2(m+1)x + m^2 + 2 = 0$ (với $m$ là tham số).
Tìm $m$ để phương trình có hai nghiệm phân biệt $x_1, x_2$ thỏa mãn
$x_1^2 + x_2^2 = 10$.
\end{de_cau}

\begin{loi_giai_cau}
Để phương trình có hai nghiệm phân biệt:
\[\Delta' = (m+1)^2 - (m^2+2) = 2m-1 > 0 \Rightarrow m > \tfrac{1}{2}.\]

Theo Vi-ét:
\[ \begin{cases} x_1+x_2 = 2(m+1) \\ x_1 x_2 = m^2+2 \end{cases} \]

\[x_1^2+x_2^2 = (x_1+x_2)^2 - 2x_1x_2 = 4(m+1)^2 - 2(m^2+2) = 2m^2+8m = 10.\]
\[m^2+4m-5=0 \Rightarrow (m+5)(m-1)=0 \Rightarrow m=1\text{ hoặc }m=-5.\]

Kết hợp $m > \frac{1}{2}$: $\boxed{m = 1}$.
\end{loi_giai_cau}
